\IfFileExists{stacks-project.cls}{%
\documentclass{stacks-project}
}{%
\documentclass{amsart}
}

% For dealing with references we use the comment environment
\usepackage{verbatim}
\newenvironment{reference}{\comment}{\endcomment}
%\newenvironment{reference}{}{}
\newenvironment{slogan}{\comment}{\endcomment}
\newenvironment{history}{\comment}{\endcomment}

% For commutative diagrams we use Xy-pic
\usepackage[all]{xy}

% We use 2cell for 2-commutative diagrams.
\xyoption{2cell}
\UseAllTwocells

% We use multicol for the list of chapters between chapters
\usepackage{multicol}

% This is generally recommended for better output
\usepackage{lmodern}
\usepackage[T1]{fontenc}

% For cross-file-references
\usepackage{xr-hyper}

% Package for hypertext links:
\usepackage{hyperref}

% For any local file, say "hello.tex" you want to link to please
% use \externaldocument[hello-]{hello}
\externaldocument[introduction-]{introduction}
\externaldocument[conventions-]{conventions}
\externaldocument[sets-]{sets}
\externaldocument[categories-]{categories}
\externaldocument[topology-]{topology}
\externaldocument[sheaves-]{sheaves}
\externaldocument[sites-]{sites}
\externaldocument[stacks-]{stacks}
\externaldocument[fields-]{fields}
\externaldocument[algebra-]{algebra}
\externaldocument[brauer-]{brauer}
\externaldocument[homology-]{homology}
\externaldocument[derived-]{derived}
\externaldocument[simplicial-]{simplicial}
\externaldocument[more-algebra-]{more-algebra}
\externaldocument[smoothing-]{smoothing}
\externaldocument[modules-]{modules}
\externaldocument[sites-modules-]{sites-modules}
\externaldocument[injectives-]{injectives}
\externaldocument[cohomology-]{cohomology}
\externaldocument[sites-cohomology-]{sites-cohomology}
\externaldocument[dga-]{dga}
\externaldocument[dpa-]{dpa}
\externaldocument[sdga-]{sdga}
\externaldocument[hypercovering-]{hypercovering}
\externaldocument[schemes-]{schemes}
\externaldocument[constructions-]{constructions}
\externaldocument[properties-]{properties}
\externaldocument[morphisms-]{morphisms}
\externaldocument[coherent-]{coherent}
\externaldocument[divisors-]{divisors}
\externaldocument[limits-]{limits}
\externaldocument[varieties-]{varieties}
\externaldocument[topologies-]{topologies}
\externaldocument[descent-]{descent}
\externaldocument[perfect-]{perfect}
\externaldocument[more-morphisms-]{more-morphisms}
\externaldocument[flat-]{flat}
\externaldocument[groupoids-]{groupoids}
\externaldocument[more-groupoids-]{more-groupoids}
\externaldocument[etale-]{etale}
\externaldocument[chow-]{chow}
\externaldocument[intersection-]{intersection}
\externaldocument[pic-]{pic}
\externaldocument[weil-]{weil}
\externaldocument[adequate-]{adequate}
\externaldocument[dualizing-]{dualizing}
\externaldocument[duality-]{duality}
\externaldocument[discriminant-]{discriminant}
\externaldocument[derham-]{derham}
\externaldocument[local-cohomology-]{local-cohomology}
\externaldocument[algebraization-]{algebraization}
\externaldocument[curves-]{curves}
\externaldocument[resolve-]{resolve}
\externaldocument[models-]{models}
\externaldocument[functors-]{functors}
\externaldocument[equiv-]{equiv}
\externaldocument[pione-]{pione}
\externaldocument[etale-cohomology-]{etale-cohomology}
\externaldocument[proetale-]{proetale}
\externaldocument[relative-cycles-]{relative-cycles}
\externaldocument[more-etale-]{more-etale}
\externaldocument[trace-]{trace}
\externaldocument[crystalline-]{crystalline}
\externaldocument[spaces-]{spaces}
\externaldocument[spaces-properties-]{spaces-properties}
\externaldocument[spaces-morphisms-]{spaces-morphisms}
\externaldocument[decent-spaces-]{decent-spaces}
\externaldocument[spaces-cohomology-]{spaces-cohomology}
\externaldocument[spaces-limits-]{spaces-limits}
\externaldocument[spaces-divisors-]{spaces-divisors}
\externaldocument[spaces-over-fields-]{spaces-over-fields}
\externaldocument[spaces-topologies-]{spaces-topologies}
\externaldocument[spaces-descent-]{spaces-descent}
\externaldocument[spaces-perfect-]{spaces-perfect}
\externaldocument[spaces-more-morphisms-]{spaces-more-morphisms}
\externaldocument[spaces-flat-]{spaces-flat}
\externaldocument[spaces-groupoids-]{spaces-groupoids}
\externaldocument[spaces-more-groupoids-]{spaces-more-groupoids}
\externaldocument[bootstrap-]{bootstrap}
\externaldocument[spaces-pushouts-]{spaces-pushouts}
\externaldocument[spaces-chow-]{spaces-chow}
\externaldocument[groupoids-quotients-]{groupoids-quotients}
\externaldocument[spaces-more-cohomology-]{spaces-more-cohomology}
\externaldocument[spaces-simplicial-]{spaces-simplicial}
\externaldocument[spaces-duality-]{spaces-duality}
\externaldocument[formal-spaces-]{formal-spaces}
\externaldocument[restricted-]{restricted}
\externaldocument[spaces-resolve-]{spaces-resolve}
\externaldocument[formal-defos-]{formal-defos}
\externaldocument[defos-]{defos}
\externaldocument[cotangent-]{cotangent}
\externaldocument[examples-defos-]{examples-defos}
\externaldocument[algebraic-]{algebraic}
\externaldocument[examples-stacks-]{examples-stacks}
\externaldocument[stacks-sheaves-]{stacks-sheaves}
\externaldocument[criteria-]{criteria}
\externaldocument[artin-]{artin}
\externaldocument[quot-]{quot}
\externaldocument[stacks-properties-]{stacks-properties}
\externaldocument[stacks-morphisms-]{stacks-morphisms}
\externaldocument[stacks-limits-]{stacks-limits}
\externaldocument[stacks-cohomology-]{stacks-cohomology}
\externaldocument[stacks-perfect-]{stacks-perfect}
\externaldocument[stacks-introduction-]{stacks-introduction}
\externaldocument[stacks-more-morphisms-]{stacks-more-morphisms}
\externaldocument[stacks-geometry-]{stacks-geometry}
\externaldocument[moduli-]{moduli}
\externaldocument[moduli-curves-]{moduli-curves}
\externaldocument[examples-]{examples}
\externaldocument[exercises-]{exercises}
\externaldocument[guide-]{guide}
\externaldocument[desirables-]{desirables}
\externaldocument[coding-]{coding}
\externaldocument[obsolete-]{obsolete}
\externaldocument[fdl-]{fdl}
\externaldocument[index-]{index}

% Theorem environments.
%
\theoremstyle{plain}
\newtheorem{theorem}[subsection]{Theorem}
\newtheorem{proposition}[subsection]{Proposition}
\newtheorem{lemma}[subsection]{Lemma}

\theoremstyle{definition}
\newtheorem{definition}[subsection]{Definition}
\newtheorem{example}[subsection]{Example}
\newtheorem{exercise}[subsection]{Exercise}
\newtheorem{situation}[subsection]{Situation}

\theoremstyle{remark}
\newtheorem{remark}[subsection]{Remark}
\newtheorem{remarks}[subsection]{Remarks}

\numberwithin{equation}{subsection}

% Macros
%
\def\lim{\mathop{\mathrm{lim}}\nolimits}
\def\colim{\mathop{\mathrm{colim}}\nolimits}
\def\Spec{\mathop{\mathrm{Spec}}}
\def\Hom{\mathop{\mathrm{Hom}}\nolimits}
\def\Ext{\mathop{\mathrm{Ext}}\nolimits}
\def\SheafHom{\mathop{\mathcal{H}\!\mathit{om}}\nolimits}
\def\SheafExt{\mathop{\mathcal{E}\!\mathit{xt}}\nolimits}
\def\Sch{\mathit{Sch}}
\def\Mor{\mathop{\mathrm{Mor}}\nolimits}
\def\Ob{\mathop{\mathrm{Ob}}\nolimits}
\def\Sh{\mathop{\mathit{Sh}}\nolimits}
\def\NL{\mathop{N\!L}\nolimits}
\def\CH{\mathop{\mathrm{CH}}\nolimits}
\def\proetale{{pro\text{-}\acute{e}tale}}
\def\etale{{\acute{e}tale}}
\def\QCoh{\mathit{QCoh}}
\def\Ker{\mathop{\mathrm{Ker}}}
\def\Im{\mathop{\mathrm{Im}}}
\def\Coker{\mathop{\mathrm{Coker}}}
\def\Coim{\mathop{\mathrm{Coim}}}

% Boxtimes
%
\DeclareMathSymbol{\boxtimes}{\mathbin}{AMSa}{"02}

%
% Macros for moduli stacks/spaces
%
\def\QCohstack{\mathcal{QC}\!\mathit{oh}}
\def\Cohstack{\mathcal{C}\!\mathit{oh}}
\def\Spacesstack{\mathcal{S}\!\mathit{paces}}
\def\Quotfunctor{\mathrm{Quot}}
\def\Hilbfunctor{\mathrm{Hilb}}
\def\Curvesstack{\mathcal{C}\!\mathit{urves}}
\def\Polarizedstack{\mathcal{P}\!\mathit{olarized}}
\def\Complexesstack{\mathcal{C}\!\mathit{omplexes}}
% \Pic is the operator that assigns to X its picard group, usage \Pic(X)
% \Picardstack_{X/B} denotes the Picard stack of X over B
% \Picardfunctor_{X/B} denotes the Picard functor of X over B
\def\Pic{\mathop{\mathrm{Pic}}\nolimits}
\def\Picardstack{\mathcal{P}\!\mathit{ic}}
\def\Picardfunctor{\mathrm{Pic}}
\def\Deformationcategory{\mathcal{D}\!\mathit{ef}}
\usepackage{cleveref}
\crefname{originalproposition}{original proposition}{original propositions}
\Crefname{originalproposition}{Original proposition}{Original propositions}
\makeatletter
\newlabel{varieties-proposition-finite-set-of-points-of-codim-1-in-affine}{{09NN}{source}{}{sourcepointer-1}{}}
\newlabel{varieties-proposition-finite-set-of-points-of-codim-1-in-affine@cref}{{[originalproposition][0][]09NN}{[1][0][]source}{}{}{}}
\makeatother
\crefname{originallemma}{original lemma}{original lemmas}
\Crefname{originallemma}{Original lemma}{Original lemmas}
\makeatletter
\newlabel{lemma-find-affine-integral}{{lemma-find-affine-integral}{source}{}{sourcepointer-2}{}}
\newlabel{lemma-find-affine-integral@cref}{{[originallemma][0][]lemma-find-affine-integral}{[1][0][]source}{}{}{}}
\makeatother
\crefname{originallemma}{original lemma}{original lemmas}
\Crefname{originallemma}{Original lemma}{Original lemmas}
\makeatletter
\newlabel{lemma-finite-flat-over-almost-dense-subscheme}{{lemma-finite-flat-over-almost-dense-subscheme}{source}{}{sourcepointer-3}{}}
\newlabel{lemma-finite-flat-over-almost-dense-subscheme@cref}{{[originallemma][0][]lemma-finite-flat-over-almost-dense-subscheme}{[1][0][]source}{}{}{}}
\makeatother
\crefname{originallemma}{original lemma}{original lemmas}
\Crefname{originallemma}{Original lemma}{Original lemmas}
\makeatletter
\newlabel{groupoids-lemma-find-invariant-affine}{{03JE}{source}{}{sourcepointer-4}{}}
\newlabel{groupoids-lemma-find-invariant-affine@cref}{{[originallemma][0][]03JE}{[1][0][]source}{}{}{}}
\makeatother
\crefname{originallemma}{original lemma}{original lemmas}
\Crefname{originallemma}{Original lemma}{Original lemmas}
\makeatletter
\newlabel{properties-lemma-ample-finite-set-in-affine}{{01ZY}{source}{}{sourcepointer-5}{}}
\newlabel{properties-lemma-ample-finite-set-in-affine@cref}{{[originallemma][0][]01ZY}{[1][0][]source}{}{}{}}
\makeatother
\crefname{originallemma}{original lemma}{original lemmas}
\Crefname{originallemma}{Original lemma}{Original lemmas}
\makeatletter
\newlabel{lemma-no-specializations-map-to-same-point}{{lemma-no-specializations-map-to-same-point}{source}{}{sourcepointer-6}{}}
\newlabel{lemma-no-specializations-map-to-same-point@cref}{{[originallemma][0][]lemma-no-specializations-map-to-same-point}{[1][0][]source}{}{}{}}
\makeatother
\crefname{originallemma}{original lemma}{original lemmas}
\Crefname{originallemma}{Original lemma}{Original lemmas}
\makeatletter
\newlabel{lemma-find-almost-invariant-function}{{lemma-find-almost-invariant-function}{source}{}{sourcepointer-7}{}}
\newlabel{lemma-find-almost-invariant-function@cref}{{[originallemma][0][]lemma-find-almost-invariant-function}{[1][0][]source}{}{}{}}
\makeatother
\crefname{originallemma}{original lemma}{original lemmas}
\Crefname{originallemma}{Original lemma}{Original lemmas}
\makeatletter
\newlabel{groupoids-lemma-determinant-trick}{{03BH}{source}{}{sourcepointer-8}{}}
\newlabel{groupoids-lemma-determinant-trick@cref}{{[originallemma][0][]03BH}{[1][0][]source}{}{}{}}
\makeatother
\crefname{originallemma}{original lemma}{original lemmas}
\Crefname{originallemma}{Original lemma}{Original lemmas}
\makeatletter
\newlabel{topology-lemma-characterize-closed-Noetherian}{{0542}{source}{}{sourcepointer-9}{}}
\newlabel{topology-lemma-characterize-closed-Noetherian@cref}{{[originallemma][0][]0542}{[1][0][]source}{}{}{}}
\makeatother
\crefname{originallemma}{original lemma}{original lemmas}
\Crefname{originallemma}{Original lemma}{Original lemmas}
\makeatletter
\newlabel{topology-lemma-constructible-stable-specialization-closed}{{0903}{source}{}{sourcepointer-10}{}}
\newlabel{topology-lemma-constructible-stable-specialization-closed@cref}{{[originallemma][0][]0903}{[1][0][]source}{}{}{}}
\makeatother
\crefname{originallemma}{original lemma}{original lemmas}
\Crefname{originallemma}{Original lemma}{Original lemmas}
\makeatletter
\newlabel{lemma-get-affine}{{lemma-get-affine}{source}{}{sourcepointer-11}{}}
\newlabel{lemma-get-affine@cref}{{[originallemma][0][]lemma-get-affine}{[1][0][]source}{}{}{}}
\makeatother
\title{Invariant affine open containing a point whose orbit has one-dimensional local rings in a finite groupoid}
\author{The Stacks Project Authors}
\date{stacks-project:more-groupoids.tex:local-snapshot-20261012; GNU FDL 1.2 or later}
\begin{document}
\maketitle
\setcounter{subsection}{1}
\noindent\textit{Editorial note.} All delivered mathematical text is the author's own, copied from the pinned Stacks Project chapter file more-groupoids.tex; no statement or proof step was rewritten, reordered or removed. The counted span is the author environment and its own proof environment. References inside the span are handed to the assembler as navigation stubs that carry the canonical tag of the cited tag only; whether each of them is load-bearing is recorded on the stub itself.\par\medskip
\noindent\textit{Allowed standard prerequisites (editorial).}\begin{itemize}
\item A finite set of points of codimension at most one in a separated locally Noetherian scheme is contained in an affine open (varieties chapter, Proposition varieties-proposition-finite-set-of-points-of-codim-1-in-affine, canonical tag 09NN). cited at proof lines 3035-3038. Puts the finite R-orbit of u into an affine open U', starting the reduction chain.
\item Descent of invariant affine opens along integral morphisms met by every orbit (more-groupoids chapter, Lemma lemma-find-affine-integral, no canonical tag in the snapshot; no-tag pointer to the source label line). cited at proof lines 3046-3047 and 3052-3053. Justifies replacing U by its reduction and later by the closure of the invariant dense open, moving u to an orbit point.
\item Existence of an invariant open W and an invariant closed W' in W with W a thickening of W', the orbit image condition on the closure of W, and finite locally free restrictions on W' (more-groupoids chapter, Lemma lemma-finite-flat-over-almost-dense-subscheme, no canonical tag in the snapshot; no-tag pointer to the source label line). cited at proof lines 3047-3049. Supplies the dense invariant open W with finite locally free structure maps used in the main construction.
\item Invariant affine open containing a finite set of points for a groupoid with finite locally free structure maps (more-groupoids chapter, Lemma lemma-find-invariant-affine, canonical tag 03JE). cited at proof lines 3060-3061. Settles the case u in W directly.
\item Any finite set of points of a quasi-affine scheme is contained in an affine open (properties chapter, Lemma properties-lemma-ample-finite-set-in-affine, canonical tag 01ZY). cited at proof lines 3062-3064. Input used to apply lemma-find-invariant-affine to the quasi-affine W.
\item Distinct points in the same orbit of a groupoid with finite s, t are never specializations of one another (more-groupoids chapter, Lemma lemma-no-specializations-map-to-same-point, no canonical tag in the snapshot; no-tag pointer to the source label line). cited at proof lines 3068-3070. Ensures a point specializing to an orbit point lies outside the orbit, so assumption (3) makes it a generic point of U.
\item Almost invariant functions: for a groupoid with s, t finite of finite presentation, points with generic orbits, an immersion into an affine scheme and an ideal I disjoint from the image with closed union, some element of I cuts out an invariant affine open containing the points (more-groupoids chapter, Lemma lemma-find-almost-invariant-function, no canonical tag in the snapshot; no-tag pointer to the source label line). cited at proof lines 3082-3086 and 3110-3113. Produces f, giving an invariant affine open inside W containing all xi\_j, and h, giving an invariant affine open inside Z containing u; these are the two pieces glued into V.
\item The norm of the pullback of a function along the source map of a groupoid is invariant (more-groupoids chapter, Lemma lemma-determinant-trick, canonical tag 03BH). cited at proof lines 3116-3117. Makes g, the norm of the pullback of h along the structure maps, an R\_W-invariant function.
\item In Noetherian topological spaces the closed subsets are exactly the constructible specialization closed subsets (topology chapter, Lemma topology-lemma-characterize-closed-Noetherian, canonical tag 0542). cited at proof lines 3130-3133. Reduces openness of the constructible set V to a generalization closedness check, carried out by the explicit chase.
\item Constructible subsets closed under specialization characterize closed subsets, general version (topology chapter, Lemma topology-lemma-constructible-stable-specialization-closed, canonical tag 0903). cited at proof lines 3133-3135. Cited alongside 0542 as the more general form of the same openness criterion.
\item Affineness criterion: a quasi-compact immersion into an affine scheme whose preimage of a principal open is affine and whose intersection with the principal closed part is closed has affine source (more-groupoids chapter, Lemma lemma-get-affine, no canonical tag in the snapshot; no-tag pointer to the source label line). cited at proof lines 3164-3166. Concludes that V is affine from the immersion into the affine scheme A\_h localized.
\end{itemize}\medskip
\ifdefined\hypertarget\hypertarget{sourcepointer-1}{}\fi
\noindent\textit{Source reference.} Original proposition 09NN: sources/stacks-chow/tags.txt, source lines 12221--12221. Cited at proof lines 3035-3038: puts the finite orbit of u into an affine open.\par\medskip
\ifdefined\hypertarget\hypertarget{sourcepointer-2}{}\fi
\noindent\textit{Source reference.} Original lemma lemma-find-affine-integral: sources/stacks-more-groupoids/pinned-original-source.tex, source lines 2887--2887. Cited at proof lines 3046 and 3052-3053: reduction to reduced U and replacement by the closure of W.\par\medskip
\ifdefined\hypertarget\hypertarget{sourcepointer-3}{}\fi
\noindent\textit{Source reference.} Original lemma lemma-finite-flat-over-almost-dense-subscheme: sources/stacks-more-groupoids/pinned-original-source.tex, source lines 2672--2672. Cited at proof lines 3047-3049: supplies the invariant dense open W with finite locally free restrictions.\par\medskip
\ifdefined\hypertarget\hypertarget{sourcepointer-4}{}\fi
\noindent\textit{Source reference.} Original lemma 03JE: sources/stacks-chow/tags.txt, source lines 4636--4636. Cited at proof lines 3060-3061: settles the case u in W.\par\medskip
\ifdefined\hypertarget\hypertarget{sourcepointer-5}{}\fi
\noindent\textit{Source reference.} Original lemma 01ZY: sources/stacks-chow/tags.txt, source lines 2613--2613. Cited at proof lines 3062-3064: input for applying 03JE to the quasi-affine W.\par\medskip
\ifdefined\hypertarget\hypertarget{sourcepointer-6}{}\fi
\noindent\textit{Source reference.} Original lemma lemma-no-specializations-map-to-same-point: sources/stacks-more-groupoids/pinned-original-source.tex, source lines 2965--2965. Cited at proof lines 3068-3070: points specializing into the orbit lie outside the orbit.\par\medskip
\ifdefined\hypertarget\hypertarget{sourcepointer-7}{}\fi
\noindent\textit{Source reference.} Original lemma lemma-find-almost-invariant-function: sources/stacks-more-groupoids/pinned-original-source.tex, source lines 2923--2923. Cited at proof lines 3082-3086 and 3110-3113: produces f and h, the two affine pieces of V.\par\medskip
\ifdefined\hypertarget\hypertarget{sourcepointer-8}{}\fi
\noindent\textit{Source reference.} Original lemma 03BH: sources/stacks-chow/tags.txt, source lines 4348--4348. Cited at proof lines 3116-3117: makes the norm function g invariant.\par\medskip
\ifdefined\hypertarget\hypertarget{sourcepointer-9}{}\fi
\noindent\textit{Source reference.} Original lemma 0542: sources/stacks-chow/tags.txt, source lines 6625--6625. Cited at proof lines 3130-3133: reduces openness of V to the generalization check.\par\medskip
\ifdefined\hypertarget\hypertarget{sourcepointer-10}{}\fi
\noindent\textit{Source reference.} Original lemma 0903: sources/stacks-chow/tags.txt, source lines 11396--11396. Cited at proof lines 3133-3135: the more general form of the same criterion.\par\medskip
\ifdefined\hypertarget\hypertarget{sourcepointer-11}{}\fi
\noindent\textit{Source reference.} Original lemma lemma-get-affine: sources/stacks-more-groupoids/pinned-original-source.tex, source lines 2987--2987. Cited at proof lines 3164-3166: concludes V is affine.\par\medskip
\begin{lemma}
\label{lemma-find-affine-codimension-1}
Let $(U, R, s, t, c)$ be a groupoid scheme. Let $u \in U$. Assume
\begin{enumerate}
\item $s, t$ are finite morphisms,
\item $U$ is separated and locally Noetherian,
\item $\dim(\mathcal{O}_{U, u'}) \leq 1$ for every point $u'$
in the orbit of $u$.
\end{enumerate}
Then $u$ is contained in an $R$-invariant affine open of $U$.
\end{lemma}

\begin{proof}
The $R$-orbit of $u$ is finite. By conditions (2) and (3) it is contained
in an affine open $U'$ of $U$, see
Varieties, Proposition
\ref{varieties-proposition-finite-set-of-points-of-codim-1-in-affine}.
Then $t(s^{-1}(U \setminus U'))$ is an $R$-invariant
closed subset of $U$ which does not contain $u$. Thus
$U \setminus t(s^{-1}(U \setminus U'))$ is an $R$-invariant open
of $U'$ containing $u$.
Replacing $U$ by this open we may assume $U$ is quasi-affine.

\medskip\noindent
By Lemma \ref{lemma-find-affine-integral} we may replace $U$ by its reduction
and assume $U$ is reduced. This means $R$-invariant subschemes
$W' \subset W \subset U$ of
Lemma \ref{lemma-finite-flat-over-almost-dense-subscheme}
are equal $W' = W$. As $U = t(s^{-1}(\overline{W}))$ some point
$u'$ of the $R$-orbit of $u$ is contained in $\overline{W}$
and by Lemma \ref{lemma-find-affine-integral}
we may replace $U$ by $\overline{W}$ and $u$ by $u'$.
Hence we may assume there is
a dense open $R$-invariant subscheme $W \subset U$ such that
the morphisms $s_W, t_W$ of the restriction $(W, R_W, s_W, t_W, c_W)$ are
finite locally free.

\medskip\noindent
If $u \in W$ then we are done by
Groupoids, Lemma \ref{groupoids-lemma-find-invariant-affine}
(because $W$ is quasi-affine so any finite set of points
of $W$ is contained in an affine open, see
Properties, Lemma \ref{properties-lemma-ample-finite-set-in-affine}).
Thus we assume $u \not \in W$ and hence none of the points of the
orbit of $u$ is in $W$. Let $\xi \in U$
be a point with a nontrivial specialization to a point $u'$ in the orbit
of $u$. Since there are no specializations among the points in the
orbit of $u$ (Lemma \ref{lemma-no-specializations-map-to-same-point})
we see that $\xi$ is not in the orbit.
By assumption (3) we see that $\xi$ is a generic point of $U$
and hence $\xi \in W$.
As $U$ is Noetherian there are finitely many of these
points $\xi_1, \ldots, \xi_m \in W$. Because $s_W, t_W$ are flat the orbit
of each $\xi_j$ consists of generic points of irreducible components
of $W$ (and hence $U$).

\medskip\noindent
Let $j : U \to \Spec(A)$ be an immersion of $U$ into an affine scheme
(this is possible as $U$ is quasi-affine). Let $J \subset A$
be an ideal such that $V(J) \cap j(W) = \emptyset$ and $V(J) \cup j(W)$
is closed. Apply Lemma \ref{lemma-find-almost-invariant-function}
to the groupoid scheme $(W, R_W, s_W, t_W, c_W)$, the morphism
$j|_W : W \to \Spec(A)$, the points $\xi_j$, and the ideal $J$
to find an $f \in J$ such that $(j|_W)^{-1}D(f)$ is an $R_W$-invariant
affine open containing $\xi_j$ for all $j$. Since $f \in J$
we see that $j^{-1}D(f) \subset W$, i.e., $j^{-1}D(f)$ is
an $R$-invariant affine open of $U$ contained in $W$
containing all $\xi_j$.

\medskip\noindent
Let $Z$ be the reduced induced closed subscheme structure on
$$
U \setminus j^{-1}D(f) = j^{-1}V(f).
$$
Then $Z$ is set theoretically
$R$-invariant (but it may not be scheme theoretically $R$-invariant).
Let $(Z, R_Z, s_Z, t_Z, c_Z)$ be the restriction of $R$ to $Z$.
Since $Z \to U$ is finite, it follows that $s_Z$ and $t_Z$ are finite.
Since $u \in Z$ the orbit of $u$ is in $Z$ and agrees with the
$R_Z$-orbit of $u$ viewed as a point of $Z$. Since
$\dim(\mathcal{O}_{U, u'}) \leq 1$ and since $\xi_j \not \in Z$
for all $j$, we see that $\dim(\mathcal{O}_{Z, u'}) \leq 0$ for
all $u'$ in the orbit of $u$. In other words, the $R_Z$-orbit of $u$
consists of generic points of irreducible components of $Z$.

\medskip\noindent
Let $I \subset A$ be an ideal such that $V(I) \cap j(U) =\emptyset$
and $V(I) \cup j(U)$ is closed. Apply
Lemma \ref{lemma-find-almost-invariant-function} to
the groupoid scheme $(Z, R_Z, s_Z, t_Z, c_Z)$, the restriction $j|_Z$,
the ideal $I$, and the point $u \in Z$ to obtain $h \in I$ such that
$j^{-1}D(h) \cap Z$ is an $R_Z$-invariant open affine containing $u$.

\medskip\noindent
Consider the $R_W$-invariant (Groupoids, Lemma
\ref{groupoids-lemma-determinant-trick}) function
$$
g = 
\text{Norm}_{s_W}(t_W^\sharp(j^\sharp(h)|_W)) \in \Gamma(W, \mathcal{O}_W)
$$
(In the following we only need the restriction of $g$ to $j^{-1}D(f)$ and
in this case the norm is along a finite locally free morphism of affines.)
We claim that
$$
V = (W_g \cap j^{-1}D(f)) \cup (j^{-1}D(h) \cap Z)
$$
is an $R$-invariant affine open of $U$ which finishes the proof of the lemma.
It is set theoretically $R$-invariant by construction. As $V$ is a
constructible set, to see that it is open it suffices to show it is
closed under generalization in $U$ (Topology, Lemma
\ref{topology-lemma-characterize-closed-Noetherian}
or the more general
Topology, Lemma
\ref{topology-lemma-constructible-stable-specialization-closed}).
Since $W_g \cap j^{-1}D(f)$ is open in $U$, it suffices to consider
a specialization $u_1 \leadsto u_2$ of $U$ with
$u_2 \in j^{-1}D(h) \cap Z$.
This means that $h$ is nonzero in $j(u_2)$ and $u_2 \in Z$.
If $u_1 \in Z$, then $j(u_1) \leadsto j(u_2)$ and since
$h$ is nonzero in $j(u_2)$ it is nonzero in $j(u_1)$ which
implies $u_1 \in V$. If $u_1 \not \in Z$ and
also not in $W_g \cap j^{-1}D(f)$, then $u_1 \in W$, $u_1 \not \in W_g$
because the complement of $Z = j^{-1}V(f)$ is contained in $W \cap j^{-1}D(f)$.
Hence there exists a point $r_1 \in R$ with $s(r_1) = u_1$
such that $h$ is zero in $t(r_1)$. Since $s$ is finite we
can find a specialization $r_1 \leadsto r_2$ with $s(r_2) = u_2$.
However, then we conclude that $h$ is zero in $u'_2 = t(r_2)$
which contradicts the fact that $j^{-1}D(h) \cap Z$
is $R$-invariant and $u_2$ is in it. Thus $V$ is open.

\medskip\noindent
Observe that $V \subset j^{-1}D(h)$ for our function $h \in I$.
Thus we obtain an immersion
$$
j' : V \longrightarrow \Spec(A_h)
$$
Let $f' \in A_h$ be the image of $f$. Then $(j')^{-1}D(f')$
is the principal open determined by $g$ in the affine
open $j^{-1}D(f)$ of $U$.
Hence $(j')^{-1}D(f)$ is affine. Finally,
$j'(V) \cap V(f') = j'(j^{-1}D(h) \cap Z)$
is closed in $\Spec(A_h/(f')) = \Spec((A/f)_h) = D(h) \cap V(f)$
by our choice of $h \in I$ and the ideal $I$. Hence we can apply
Lemma \ref{lemma-get-affine}
to conclude that $V$ is affine as claimed above.
\end{proof}

\end{document}
